%% ==========================================================================
%%  SECTION 10 --- CONCLUSION
%% ==========================================================================
\section{Conclusion}
\label{sec:conclusion}

Privacy-preserving recommendation is not one problem but four, corresponding to the stages of the
pipeline in Figure~\ref{fig:pipeline}. Reading the literature through that lens is what this
survey has tried to make possible, and it explains why systems that appear to be alternatives
frequently are not.

Two of those stages are in good shape. Private training has advanced to the point where
cryptographic privacy costs compute and communication rather than recommendation quality:
\nudgeplain{} reports nDCG@20 on par with non-private matrix factorisation~\cite{nudge2026}.
Private retrieval has developed rapidly and independently, from secure $k$-nearest-neighbour
search through web-scale private search to index-based sublinear schemes. The delivery stage is
served by a mature private-information-retrieval literature that is, with one exception, not
connected to any recommender.

That exception, \pirsonaplain{}~\cite{pirsona2021}, is also the reason the gap is worth naming:
it demonstrates that collection, training and delivery can be addressed together, and its
training core has since been superseded while the retrieval literature it would now draw on did
not yet exist. The current state of the art in training, meanwhile, states the remaining stage as
a non-goal~\cite{nudge2026}.

Two caveats bound what this survey establishes. It reports each system in the terms its authors
used and declines to rank them, because Section~\ref{sec:eval} argues that cross-paper comparison
in this field is unsound. And its claims about individual works are bounded by what we read: for
papers we could not obtain, we have said so at the point of use rather than paraphrasing them.

The problems collected in Section~\ref{sec:open} --- malicious security at scale, weakening the
non-collusion assumption, connecting delivery to training, and the leakage that survives any
protocol because it is carried by the recommendation itself --- are the ones the surveyed authors
themselves identify. The last of these is the most fundamental, since no improvement in protocol
design addresses it.
